% =====================================================================
%  Anteproyecto: WikiLLM como base de conocimiento incremental
%  Universidad Icesi - Maestria en Inteligencia Artificial Aplicada
%  Formato: plantilla de presentacion de documentos, Facultad Barberi.
%  Compilar con: pdflatex anteproyecto.tex  (dos veces)
% =====================================================================
\documentclass[11pt,letterpaper]{article}

\usepackage[utf8]{inputenc}
\usepackage[T1]{fontenc}
\usepackage[spanish,es-tabla,es-nodecimaldot]{babel}
\usepackage{mathptmx}                 % Times New Roman
\usepackage[scaled=0.92]{helvet}
\usepackage{courier}

\usepackage[letterpaper,top=2.54cm,bottom=2.54cm,left=2.54cm,right=2.54cm,
            headheight=14pt,headsep=18pt]{geometry}

\usepackage{fancyhdr}
\usepackage{titlesec}
\usepackage{booktabs}
\usepackage{array}
\usepackage{ragged2e}
\usepackage{caption}
\usepackage{enumitem}
\usepackage[hidelinks]{hyperref}
\usepackage{xurl}

% ---------------------------------------------------------------- estilo
\pagestyle{fancy}
\fancyhf{}
\fancyhead[L]{\footnotesize Universidad Icesi, Facultad Barberi de Ingeniería Diseño y Ciencias Aplicadas}
\fancyhead[R]{\footnotesize\thepage}
\renewcommand{\headrulewidth}{0.4pt}

\titleformat{\section}[block]{\centering\normalsize\scshape}{\Roman{section}.\quad}{0pt}{}
\titlespacing*{\section}{0pt}{10pt}{4pt}
\titleformat{\subsection}[block]{\normalsize\itshape}{\Alph{subsection}.\quad}{0pt}{}
\titlespacing*{\subsection}{\parindent}{5pt}{2pt}

\captionsetup[table]{labelsep=newline,justification=centering,font=sc,labelfont=sc,skip=2pt}
\renewcommand{\thetable}{\Roman{table}}
\renewcommand{\arraystretch}{1.15}

\setlength{\parindent}{1.3em}
\setlength{\parskip}{0pt}
\setlist{itemsep=1pt,topsep=3pt}

\newcolumntype{L}[1]{>{\RaggedRight\arraybackslash}p{#1}}
\newcommand{\refhang}[1]{\par\noindent\hangindent=1.5em\hangafter=1 #1\par}

% nota al pie sin marca, como en la plantilla
\newcommand{\notasinmarca}[1]{%
  \begingroup\renewcommand{\thefootnote}{}\footnote{#1}\addtocounter{footnote}{-1}\endgroup}

% ================================================================ inicio
\begin{document}

\thispagestyle{fancy}

\begin{center}
{\Large Anteproyecto\par}
\vspace{4pt}
{\LARGE WikiLLM: compilación incremental de conocimiento\\[2pt]
para optimizar el contexto de modelos de lenguaje y personas\par}
\vspace{8pt}
{\ttfamily\small Juan Diego Ramírez\quad Juan David Cruz García\par}
{\ttfamily\small juan.ramirezzuluaga@u.icesi.edu.co, juan.cruz6@u.icesi.edu.co\par}
\end{center}

\vspace{4pt}

\notasinmarca{Anteproyecto de la Maestría en Inteligencia Artificial Aplicada, Universidad Icesi,
presentado para revisión del tutor. Repositorio:
\url{https://github.com/1KVueltasAlCampo/llm-knowledge-management-wiki-agent}.}

\noindent\textbf{Resumen —} Los sistemas de Recuperación Aumentada por Generación (RAG) responden
cada consulta recuperando fragmentos y sintetizándolos de nuevo, sin conservar lo aprendido. En
el patrón WikiLLM, en cambio, un agente compila las fuentes, una a una, en páginas markdown
enlazadas que crecen y se corrigen con cada ingesta. Este anteproyecto plantea estudiar el WikiLLM como
herramienta para construir una base de conocimiento incremental que reduzca el contexto que un
modelo y una persona necesitan leer para responder correctamente. El estudio usa literatura abierta
de inteligencia artificial publicada en arXiv después del corte de entrenamiento del modelo. Tras cada ingesta se guarda una copia del wiki y
se mide, frente a una línea base RAG sobre el mismo corpus, la calidad de las respuestas, la
cantidad de contexto consumido, el esfuerzo de lectura humana y el costo. Se espera aportar
un protocolo de evaluación abierto y un primer conjunto de mediciones sobre cuándo conviene
compilar el conocimiento en lugar de recuperarlo.

\vspace{4pt}
\noindent\textbf{Palabras clave —} Base de conocimiento incremental, ingeniería de contexto,
modelos de lenguaje, recuperación aumentada por generación, WikiLLM.

% ====================================================================
\section{Planteamiento del problema}

Buscar información consume tiempo de trabajo: en un estudio con 716 empleados públicos, el
22,3\,\% dedicaba cerca de media jornada semanal a buscar la información que necesitaba
(Nakash y Bouhnik, 2024). Los asistentes basados en RAG (Lewis et al., 2020) se adoptan para reducir ese costo:
ante cada pregunta recuperan fragmentos por similitud semántica y el modelo redacta una respuesta
a partir de ellos.

Este diseño tiene tres límites. Primero, la síntesis se repite completa en cada consulta y no
se conserva (Karpathy, 2026). Segundo, el modelo recibe fragmentos sueltos, sin fecha ni señales
sobre la calidad de su fuente, y debe reconciliar en el momento datos obsoletos o
contradictorios. Tercero, la persona que verifica la respuesta solo ve esos mismos fragmentos.
En los tres casos el recurso escaso es el \textit{contexto}: los tokens que el modelo procesa
en cada llamada, con su costo y su ruido, y el tiempo de lectura de la persona. Un contexto más
grande no garantiza mejores respuestas; importa entregar la información pertinente en el menor
texto posible.

El patrón WikiLLM (Karpathy, 2026) mueve la síntesis desde la consulta hacia la ingesta: un
agente lee cada fuente nueva, actualiza las páginas afectadas de un wiki en markdown y registra
las contradicciones. Las consultas se responden sobre páginas ya redactadas, enlazadas y
fechadas, que también puede leer una persona. La Tabla~\ref{tab:comparacion}
resume la diferencia.

\begin{table}[!ht]
\centering
\caption{Comparación entre RAG y el patrón WikiLLM}
\label{tab:comparacion}
\small
\begin{tabular}{L{4.5cm} L{4.6cm} L{5.0cm}}
\toprule
 & \textbf{RAG (línea base)} & \textbf{WikiLLM (propuesta)} \\
\midrule
Momento de la síntesis & En cada consulta & Una vez, al ingerir cada fuente \\
Qué se conserva & Fragmentos y vectores & Páginas redactadas y enlazadas \\
Contexto para el modelo & Fragmentos crudos por similitud & Páginas consolidadas, con fuente \\
Contexto para la persona & Fragmentos sueltos & Wiki navegable y versionado \\
Manejo de contradicciones & Implícito, al responder & Explícito, marcado y fechado \\
Falla característica & Recuperar un fragmento erróneo & Conservar un error de síntesis \\
\bottomrule
\end{tabular}
\end{table}

El estudio usa literatura abierta de inteligencia artificial en arXiv, un caso exigente del
problema: en febrero de 2026 llegaron más de 20.500 preprints de computación sin revisión por
pares (arXiv, 2026), y el texto generado por modelos es casi seis veces más frecuente en los
artículos que no son revisiones (Elazar y Antoniak, 2026). Al ser público, el corpus permite
replicar el experimento.

% ====================================================================
\section{Antecedentes y vacío de investigación}

La revisión del estado del arte identificó veinte trabajos de 2026 sobre el patrón, todos
preprints, cuyos resultados frente a RAG no coinciden. Wright (2026) encuentra un empate en
hechos puntuales y una ventaja del wiki en preguntas que cruzan unidades de un curso. En cambio,
Zhang (2026) no observa ninguna variante del wiki que supere a RAG, y Huerta (2026a) reporta que
compilar sin control pierde información en el 53--60\,\% de los casos, aunque una o dos rondas
de refinamiento recuperan cerca del 80\,\% de la calidad perdida.

El argumento económico, el más ligado al contexto, también está en disputa. Wen y Ku (2026) estiman un 84,6\,\% menos de tokens acumulados que RAG, mientras que Cochran
(2026a) mide que el wiki consume unas 21 veces más tokens por consulta. Cochran (2026b) muestra
que la diferencia depende de cómo se accede al wiki: con acceso progresivo a las páginas, el
costo baja entre un tercio y más de la mitad sin pérdida de calidad. Por último, la organización
del wiki tiende a deteriorarse a medida que crece (Zhou et al., 2026).

Las comparaciones existentes evalúan un wiki ya construido, con muestras de 10 a 59 preguntas,
y reportan el costo como un valor aislado. Ningún trabajo mide, tras cada ingesta, cuánto contexto necesitan un modelo y una persona
para responder correctamente, ni si esa cantidad mejora o empeora a medida que el wiki crece y
recibe fuentes de calidad desigual.

% ====================================================================
\section{Pregunta, hipótesis y objetivos}

\noindent\textbf{Pregunta de investigación.} \textit{¿En qué medida un WikiLLM mantenido de forma
incremental reduce el contexto que un modelo de lenguaje y una persona necesitan para responder
correctamente, y cómo evoluciona esa ventaja frente a un sistema RAG a medida que se ingieren
nuevas fuentes?}

\begin{itemize}[leftmargin=2.4em]
\item[H1.] A partir de cierto número de fuentes ingeridas, el WikiLLM obtiene una exactitud igual
o mayor que RAG con menos tokens de contexto por respuesta correcta, y una persona encuentra la
respuesta en menos tiempo.
\item[H2.] (Riesgo.) Compilar puede convertir afirmaciones débiles en prosa segura y conservar
errores durante varias ingestas, lo que reduce o anula esa ventaja; si ocurre, se reportará.
\end{itemize}

\noindent\textbf{Objetivo general.} Evaluar el patrón WikiLLM como herramienta para construir una
base de conocimiento incremental que optimice el contexto disponible para modelos de lenguaje y
personas, en comparación con RAG.

\textbf{Objetivos específicos:} (1)~construir un corpus versionado de 30 a 50 preprints y
preguntas de control con respuesta verificable; (2)~implementar un agente WikiLLM que conserve la
fuente de cada afirmación y someta sus cambios a revisión humana; (3)~implementar una línea base
RAG reproducible sobre el mismo corpus; y (4)~medir tras cada ingesta la calidad, el contexto
consumido, el esfuerzo humano y el costo de ambos sistemas.

% ====================================================================
\section{Metodología}

\subsection{Corpus}
Se selecciona una subárea de inteligencia artificial con tablas de clasificación activas y una
ventana temporal posterior al corte de entrenamiento del modelo. Así, las afirmaciones son
numéricas, fechadas y verificables, y el modelo no puede responder de memoria. Cada afirmación conserva su fuente y señales
de respaldo: disponibilidad de código, estado de publicación y versión del preprint.

\subsection{Sistema}
El agente se implementa en Python con \texttt{litellm}, lo que permite cambiar de proveedor de
modelo sin modificar el código y registrar los tokens de cada llamada. Ejecuta tres operaciones:
\textit{ingesta} (leer una fuente y actualizar las páginas afectadas), \textit{consulta}
(responder con citas a páginas del wiki) y \textit{revisión} o \textit{lint} (detectar
contradicciones, datos obsoletos y páginas sin enlaces). El agente no escribe directamente: entrega
sus cambios como \textit{Pull Request} en Git y un revisor los aprueba en Obsidian, lo que deja
registrado el costo de la supervisión humana.

\subsection{Protocolo experimental}
Las fuentes se ordenan por fecha y se ingieren de a una. Tras cada ingesta se guarda una copia
versionada del wiki, y sobre cada copia se responde el mismo conjunto de sesenta preguntas de
control. Los dos autores fijan por separado cada respuesta correcta y se reporta su acuerdo. RAG usa el
mismo corpus, modelo y preguntas. Se varían la frecuencia de la revisión automática y el grado de
intervención humana, y se introducen a propósito errores y fuentes de baja calidad.

\subsection{Métricas}
En la Tabla~\ref{tab:metricas}, las dos primeras métricas responden a la pregunta; las demás
controlan que la reducción de contexto no cueste calidad.

\begin{table}[!ht]
\centering
\caption{Métricas de evaluación}
\label{tab:metricas}
\small
\begin{tabular}{L{4.2cm} L{9.9cm}}
\toprule
\textbf{Métrica} & \textbf{Qué mide} \\
\midrule
Eficiencia de contexto & Respuestas correctas por cada mil tokens de contexto del modelo. \\
Esfuerzo humano & Palabras y minutos de lectura para hallar la respuesta (15 preguntas). \\
Cobertura & Proporción de preguntas de control respondidas correctamente. \\
Fundamentación & Proporción de afirmaciones rastreables hasta su fuente. \\
Contradicciones & Si el wiki se corrige cuando una fuente nueva supera un dato. \\
Persistencia del error & Ingestas que sobrevive un error introducido a propósito. \\
Costo & Tokens y tiempo por ingesta y por consulta frente a RAG. \\
\bottomrule
\end{tabular}
\end{table}

Se usan respuestas fijadas por los autores y no solo un modelo como juez, cuyo acuerdo es bajo
(Cochran, 2026a) y que se sesga con contenido posterior a su corte (Huerta, 2026b).

% ====================================================================
\section{Resultados esperados y contribución}

Se espera obtener una curva por métrica a lo largo de las ingestas para ambos sistemas y el
punto de equilibrio: el número de consultas por fuente a partir del cual compilar resulta más
barato que recuperar. Un resultado desfavorable al WikiLLM también es válido, porque delimita
cuándo conviene cada arquitectura. La contribución es un protocolo de evaluación del
mantenimiento incremental que incluye el contexto consumido por el modelo y por la persona, un
primer conjunto de mediciones y una implementación abierta bajo licencia MIT.

% ====================================================================
\section{Plan de trabajo}

La Tabla~\ref{tab:fases} cuenta las semanas desde la aprobación; las fechas se acordarán con el tutor.

\begin{table}[!ht]
\centering
\caption{Fases y entregables}
\label{tab:fases}
\small
\begin{tabular}{L{3.4cm} L{1.5cm} L{9.1cm}}
\toprule
\textbf{Fase} & \textbf{Semanas} & \textbf{Entregable} \\
\midrule
F0 Delimitación & 1 & Subárea, ventana temporal y criterio de inclusión. \\
F1 Corpus & 2--3 & Corpus versionado con metadatos de calidad. \\
F2 Protocolo & 4--5 & Preguntas de control y script de evaluación. \\
F3 Agente WikiLLM & 6--7 & Ingesta, consulta y revisión con aprobación humana. \\
F4 Línea base & 8--9 & RAG reproducible y registro de costos. \\
F5 Experimento & 10--11 & Ingestas, variantes de revisión e inyección de errores. \\
F6 Cierre & 12--13 & Informe, guía de reproducción y sustentación. \\
\bottomrule
\end{tabular}
\end{table}

Ante retrasos se recortan, en orden, las variantes de intervención humana, las de revisión y el
número de fuentes; la línea base y la inyección de errores se mantienen.

% ====================================================================
\section{Alcance, riesgos y consideraciones éticas}

El trabajo cubre un solo dominio y de 30 a 50 fuentes; su extensión a bases corporativas queda
como trabajo futuro. Que los autores construyan las preguntas es un sesgo que se reduce,
sin eliminarse, al responderlas por separado.

En seguridad, una fuente puede traer instrucciones para el agente; por eso las fuentes se tratan
como datos y toda escritura pasa por revisión humana. En sesgo, se mide si el sistema presenta afirmaciones débiles como seguras. En privacidad, el corpus es
público y el sistema admite un modelo local.

% ====================================================================
\section*{Referencias}

\footnotesize\sloppy

\refhang{arXiv. (2026). \textit{Monthly submissions}. \url{https://arxiv.org/stats/monthly_submissions}}

\refhang{Cochran, T. O. (2026a). \textit{Single-round vector RAG vs an LLM-compiled wiki: A preregistered comparison on a small multi-domain research corpus} (arXiv:2605.18490). arXiv. \url{https://arxiv.org/abs/2605.18490}}

\refhang{Cochran, T. O. (2026b). \textit{Progressive disclosure for LLM-maintained wiki knowledge bases: A preregistered ablation} (arXiv:2607.04576). arXiv. \url{https://arxiv.org/abs/2607.04576}}

\refhang{Elazar, Y., y Antoniak, M. (2026). \textit{LLM-generated or human-written? Comparing review and non-review papers on arXiv} (arXiv:2601.17036). arXiv. \url{https://arxiv.org/abs/2601.17036}}

\refhang{Huerta, J. M. (2026a). \textit{WiCER: Wiki-memory compile, evaluate, refine. Iterative knowledge compilation for LLM wiki systems} (arXiv:2605.07068). arXiv. \url{https://arxiv.org/abs/2605.07068}}

\refhang{Huerta, J. M. (2026b). \textit{Streaming knowledge compilation: Proactive materiality-scored pinning for time-evolving LLM wikis} (arXiv:2606.09877). arXiv. \url{https://arxiv.org/abs/2606.09877}}

\refhang{Karpathy, A. (2026, 4 de abril). \textit{LLM wiki: A pattern for building personal knowledge bases using LLMs} [Gist]. GitHub. \url{https://gist.github.com/karpathy/442a6bf555914893e9891c11519de94f}}

\refhang{Lewis, P., Perez, E., Piktus, A., Petroni, F., Karpukhin, V., Goyal, N., et al. (2020). Retrieval-augmented generation for knowledge-intensive NLP tasks. En \textit{Advances in Neural Information Processing Systems 33 (NeurIPS 2020)}. \url{https://arxiv.org/abs/2005.11401}}

\refhang{Nakash, M., y Bouhnik, D. (2024). How much time does the workforce spend searching for information in the ``new normal''? En \textit{iConference 2024 Proceedings}. \url{https://www.ideals.illinois.edu/items/129980}}

\refhang{Wen, S., y Ku, B. (2026). \textit{Knowledge compounding: An empirical economic analysis of self-evolving knowledge wikis under the agentic ROI framework} (arXiv:2604.11243). arXiv. \url{https://arxiv.org/abs/2604.11243}}

\refhang{Wright, B. (2026). \textit{Potential for enhanced learning in machine learning classes by using wiki LLM indexing} (arXiv:2609.25303). arXiv. \url{https://arxiv.org/abs/2609.25303}}

\refhang{Zhang, B. (2026). \textit{Does the wiki grow smarter? A pilot evaluation of LLM-maintained knowledge bases in an exception-dense decision-support task} [Preprint]. SciErXiv. \url{https://doi.org/10.20517/scierxiv202605.0291.v1}}

\refhang{Zhou, S., Yu, S., Wei, H., Wu, J., Ouyang, S., Jiao, Y., et al. (2026). \textit{Filesystem-based memory for LLM agents: Organization, evolution, and sustainability} (arXiv:2607.26637). arXiv. \url{https://arxiv.org/abs/2607.26637}}

\normalsize

\end{document}
